% =============================================================
%  01_reference.tex — Quick Reference dettagliata
% =============================================================
\chapter*{Quick Reference}
\addcontentsline{toc}{chapter}{Quick Reference}
\markboth{Quick Reference}{Quick Reference}

% --- Box per le macro-aree della reference ---
% Usa gli stessi colori semantici del cheatsheet
% Testo interno: grigio e più piccolo per non competere col codice
% UNBREAKABLE per evitare rotture interne
\newtcolorbox{refbox}[2][]{%
  enhanced,
  colback=#2!4,
  colframe=#2!70!black,
  fonttitle=\bfseries\sffamily\normalsize,
  coltitle=white,
  title={#1},
  boxrule=0.8pt,
  arc=4pt,
  left=12pt, right=12pt, top=10pt, bottom=10pt,
  toptitle=4pt, bottomtitle=4pt,
  before skip=14pt plus 4pt,
  after skip=14pt plus 4pt,
  fontupper=\small\color{black!75},  % testo grigio, più piccolo
  parskip=6pt                         % spazio tra paragrafi interni
}

% Box per note/avvertimenti interni alle sezioni
\newtcolorbox{refnote}[1][]{%
  enhanced,
  colback=black!3,
  colframe=black!30,
  fonttitle=\bfseries\small,
  coltitle=black!70,
  title={#1},
  boxrule=0.4pt,
  arc=3pt,
  left=8pt, right=8pt, top=6pt, bottom=6pt,
  before skip=10pt,
  after skip=10pt,
  fontupper=\small\color{black!65}   % testo ancora più leggero
}

% Comando per sottotitoli interni ai box (evita break dopo)
\newcommand{\reftitle}[1]{\textbf{#1}\nopagebreak\par\nopagebreak}

\noindent
Questa sezione descrive i costrutti fondamentali di Common Lisp
in forma concisa. Per ogni argomento: sintassi, spiegazione e
almeno un esempio.

\vspace{16pt}

\begin{tcolorbox}[
  enhanced,
  colback=black!2,
  colframe=black!25,
  boxrule=0.4pt,
  arc=3pt,
  left=12pt, right=12pt, top=10pt, bottom=10pt
]
\centering
\textbf{Legenda colori}\\[8pt]
\textcolor{csBlue}{\rule{1.2em}{1.2em}}\enspace Basi e stile\qquad
\textcolor{csOrange}{\rule{1.2em}{1.2em}}\enspace Controllo del flusso\qquad
\textcolor{csGreen}{\rule{1.2em}{1.2em}}\enspace Liste e sequenze\\[6pt]
\textcolor{csTeal}{\rule{1.2em}{1.2em}}\enspace Strutture dati e I/O\qquad
\textcolor{csPurple}{\rule{1.2em}{1.2em}}\enspace Astrazione e OOP
\end{tcolorbox}

\vspace{16pt}

% =============================================================
% MACRO-AREA 1: BASI DEL LINGUAGGIO
% =============================================================
\addcontentsline{toc}{section}{Basi del linguaggio}

% --- REPL e S-expression ---
\begin{refbox}[REPL e S-expression]{csBlue}
\label{sec:repl}

Common Lisp usa la notazione \textbf{S-expression} (prefix):
il primo elemento di ogni lista è l'operatore, i successivi
sono gli argomenti.

\begin{lstlisting}
(+ 1 2 3)        ; -> 6
(* 2 (+ 3 4))    ; -> 14
\end{lstlisting}

I commenti si scrivono con \fn{;} (riga singola) o
\fn{\#|~...|~\#} (multiriga):

\begin{lstlisting}
; commento su una riga
#| commento
   su piu' righe |#
\end{lstlisting}

Il REPL legge, valuta e stampa ogni espressione.
Le variabili \fn{*}, \fn{**}, \fn{***} contengono gli ultimi
tre risultati restituiti.

\end{refbox}

% --- Numeri ---
\begin{refbox}[Numeri]{csBlue}
\label{sec:numeri}

Common Lisp distingue interi, razionali, float e complessi:

\begin{lstlisting}
42      -7       ; interi
3.14    1.5e10   ; float
1/3     22/7     ; razionali esatti
#C(1 2)          ; complesso 1+2i
\end{lstlisting}

\textbf{Operatori aritmetici:}\nopagebreak

\begin{lstlisting}
(+ a b)   (- a b)   (* a b)   (/ a b)
(mod 10 3)     ; -> 1  (segno del divisore)
(rem 10 3)     ; -> 1  (segno del dividendo)
(expt 2 10)    ; -> 1024
(sqrt 9)       ; -> 3.0
(abs -5)       ; -> 5
(1+ x) (1- x)  ; incrementa / decrementa
(incf x)       ; setf di (1+ x)
(decf x)       ; setf di (1- x)
(floor 7/2)    ; -> 3  (verso -inf)
(ceiling 7/2)  ; -> 4  (verso +inf)
(truncate 7/2) ; -> 3  (verso zero)
(round 2.5)    ; -> 2  (round-to-even)
(max 3 1 4)    ; -> 4
(min 3 1 4)    ; -> 1
\end{lstlisting}

\textbf{Confronti numerici} (usare \fn{=}, non \fn{eq}/\fn{equal}):\nopagebreak

\begin{lstlisting}
(= a b)  (/= a b)  (< a b)  (> a b)  (<= a b)  (>= a b)
\end{lstlisting}

\end{refbox}

% --- Booleani ---
\begin{refbox}[Booleani]{csBlue}
\label{sec:booleani}

\fn{NIL} è l'unico valore falso; qualsiasi altra cosa è vera.
\fn{NIL} è anche la lista vuota.

\begin{lstlisting}
T    ; vero
NIL  ; falso e lista vuota
\end{lstlisting}

\begin{lstlisting}
(and T NIL)    ; -> NIL
(or  NIL 42)   ; -> 42  (primo valore non-NIL)
(not NIL)      ; -> T
\end{lstlisting}

\fn{and} e \fn{or} usano \textit{short-circuit evaluation} e
restituiscono il valore che ha determinato il risultato,
non necessariamente \fn{T} o \fn{NIL}.

\end{refbox}

% --- Simboli, Stringhe, Caratteri ---
\begin{refbox}[Simboli, stringhe e caratteri]{csBlue}
\label{sec:simboli}

\textbf{Simboli:} identificatori unici — lo stesso nome riferisce sempre lo stesso oggetto.
La virgoletta singola \fn{'} previene la valutazione:\nopagebreak

\begin{lstlisting}
'foo          ; simbolo FOO
(eq 'a 'a)   ; -> T  (stessa identita')
\end{lstlisting}

\textbf{Stringhe:}\nopagebreak

\begin{lstlisting}
"hello"
(length "hello")           ; -> 5
(char "hello" 0)           ; -> #\h
(subseq "hello" 1 3)       ; -> "el"
(string-upcase "ciao")     ; -> "CIAO"
(string-downcase "CIAO")   ; -> "ciao"
(concatenate 'string "foo" "bar") ; -> "foobar"
(string= "abc" "abc")      ; -> T  (case-sensitive)
(string-equal "ABC" "abc") ; -> T  (case-insensitive)
(parse-integer "42")       ; -> 42
\end{lstlisting}

\textbf{Caratteri:}\nopagebreak

\begin{lstlisting}
#\A  #\space  #\newline  #\tab
(char-code #\A)          ; -> 65
(code-char 65)           ; -> #\A
(char-upcase #\a)        ; -> #\A
(alpha-char-p #\a)       ; -> T
(digit-char-p #\5)       ; -> 5  (o NIL)
\end{lstlisting}

\end{refbox}

% --- Variabili ---
\begin{refbox}[Variabili e binding]{csBlue}
\label{sec:variabili}

\textbf{Globali:}\nopagebreak
\begin{lstlisting}
(defparameter *contatore* 0)  ; globale, sovrascritta al reload
(defvar *debug* nil)          ; globale, non sovrascritta se gia' legata
(defconstant +pi-greco+ 3.14159265) ; costante
\end{lstlisting}

Convenzione: variabili globali dinamiche tra \fn{*...*} (earmuffs);
costanti tra \fn{+...+}.

\medskip
\textbf{Binding locale — \fn{let}:} tutte le variabili sono
inizializzate in parallelo, non possono riferirsi l'una all'altra:\nopagebreak
\begin{lstlisting}
(let ((x 10)
      (y 20))
  (+ x y))   ; -> 30
\end{lstlisting}

\textbf{Binding sequenziale — \fn{let*}:} ogni variabile può
usare quelle definite prima:\nopagebreak
\begin{lstlisting}
(let* ((r 5)
       (area (* pi r r)))
  area)       ; -> 78.539...
\end{lstlisting}

\textbf{Assegnamento:}\nopagebreak
\begin{lstlisting}
(setf *contatore* (1+ *contatore*))
(incf *contatore*)      ; equivalente piu' conciso
(decf *contatore* 3)    ; decrementa di 3
\end{lstlisting}

\end{refbox}

% --- Funzioni ---
\begin{refbox}[Definizione di funzioni]{csBlue}
\label{sec:funzioni}

\begin{lstlisting}
(defun quadrato (x)
  "Restituisce il quadrato di X."  ; docstring (opzionale)
  (* x x))

(quadrato 5)   ; -> 25
\end{lstlisting}

Il valore di ritorno è quello dell'ultima espressione.
\fn{return-from} permette un'uscita anticipata:

\begin{lstlisting}
(defun cerca (item lst)
  (dolist (x lst nil)          ; nil = valore se lista esaurita
    (when (equal x item)
      (return-from cerca t))))
\end{lstlisting}

\textbf{Argomenti opzionali, rest, keyword:}\nopagebreak
\begin{lstlisting}
; &optional con valore di default
(defun saluta (&optional (nome "mondo"))
  (format t "Ciao, ~A!~%" nome))

; &rest raccoglie argomenti variabili in una lista
(defun somma (&rest numeri)
  (reduce #'+ numeri :initial-value 0))

; &key usa argomenti nominati
(defun crea-rettangolo (&key (w 1) (h 1))
  (* w h))
(crea-rettangolo :w 3 :h 4)   ; -> 12
\end{lstlisting}

\textbf{Funzioni locali} con \fn{flet} e \fn{labels}:\nopagebreak
\begin{lstlisting}
; flet: funzioni locali non ricorsive
(flet ((doppio (x) (* 2 x)))
  (doppio 5))    ; -> 10

; labels: funzioni locali, anche ricorsive e mutuamente ricorsive
(labels ((pari-p (n) (if (= n 0) t (dispari-p (1- n))))
         (dispari-p (n) (if (= n 0) nil (pari-p (1- n)))))
  (pari-p 4))    ; -> T
\end{lstlisting}

\textbf{\fn{progn}, \fn{prog1}, \fn{prog2}:} valutano sequenze di espressioni:\nopagebreak
\begin{lstlisting}
(progn e1 e2 e3)   ; ritorna e3
(prog1 e1 e2 e3)   ; ritorna e1
(prog2 e1 e2 e3)   ; ritorna e2
\end{lstlisting}

\end{refbox}

% =============================================================
% MACRO-AREA 2: CONTROLLO DEL FLUSSO
% =============================================================
\addcontentsline{toc}{section}{Controllo del flusso}

% --- Condizioni ---
\begin{refbox}[Condizioni]{csOrange}
\label{sec:condizioni}

\begin{lstlisting}
; if: ramo singolo (else opzionale)
(if (> x 0) 'positivo 'non-positivo)

; cond: rami multipli, t e' il default
(cond
  ((< x 0) 'negativo)
  ((= x 0) 'zero)
  (t        'positivo))

; when/unless: no ramo else, corpo multiplo
(when (> x 0)
  (print "positivo")
  (+ x 1))

(unless (null lst)
  (car lst))

; case: switch su valore (eql), otherwise = default
(case giorno
  ((1) "lunedi")
  ((2) "martedi")
  (otherwise "altro"))

; and/or come condizioni: ritornano il valore determinante
(and (> x 0) x)      ; -> x se positivo, NIL altrimenti
(or  val-a val-b)    ; -> val-a se non-NIL, altrimenti val-b
\end{lstlisting}

\end{refbox}

% --- Iterazione ---
\begin{refbox}[Iterazione]{csOrange}
\label{sec:iterazione}

\textbf{\fn{dotimes} e \fn{dolist}:}\nopagebreak
\begin{lstlisting}
(dotimes (i 5)          ; i da 0 a 4
  (print i))

(dolist (x '(a b c) 'fatto)  ; terzo arg = valore di ritorno
  (print x))
\end{lstlisting}

\textbf{\fn{loop} — forme comuni:}\nopagebreak
\begin{lstlisting}
; iterazione su lista
(loop for x in '(1 2 3)
      collect (* x x))          ; -> (1 4 9)

; iterazione numerica
(loop for i from 1 to 5
      sum i)                     ; -> 15

; passo personalizzato
(loop for i from 0 below 10 by 2
      collect i)                 ; -> (0 2 4 6 8)

; filtro
(loop for x in '(1 2 3 4 5)
      when (evenp x)
      collect x)                 ; -> (2 4)

; con accumulatore esplicito
(loop with acc = 0
      for x in '(1 2 3 4)
      do (incf acc x)
      finally (return acc))      ; -> 10

; loop infinito con uscita esplicita
(loop
  (when (done-p stato) (return risultato))
  (setf stato (aggiorna stato)))
\end{lstlisting}

\end{refbox}

% --- Ricorsione ---
\begin{refbox}[Ricorsione]{csOrange}
\label{sec:ricorsione}

\textbf{Schema base su lista:}\nopagebreak
\begin{lstlisting}
(defun lunghezza (lst)
  (if (null lst)
      0                                ; caso base
      (1+ (lunghezza (cdr lst)))))     ; passo ricorsivo
\end{lstlisting}

\textbf{Tail recursion con accumulatore} (evita stack overflow):\nopagebreak
\begin{lstlisting}
(defun lunghezza (lst)
  (labels ((iter (lst acc)
             (if (null lst)
                 acc
                 (iter (cdr lst) (1+ acc)))))
    (iter lst 0)))
\end{lstlisting}

Una chiamata è \textit{tail-recursive} se è l'ultima operazione
prima del ritorno. Allegro CL può ottimizzare la tail recursion (TCO),
eliminando i frame dallo stack — ma questo non è garantito
dallo standard ANSI.

\begin{refnote}[Schema ricorsione su lista]
\begin{tabular}{@{}ll@{}}
\textbf{Caso base:} & \fn{(null lst)} \\
\textbf{Caso ricorsivo:} & opera su \fn{(car lst)}, continua su \fn{(cdr lst)} \\
\textbf{Con accumulatore:} & aggiunge parametro \fn{acc}, inizializzato nel wrapper \\
\end{tabular}
\end{refnote}

\end{refbox}

% =============================================================
% MACRO-AREA 3: LISTE E SEQUENZE
% =============================================================
\addcontentsline{toc}{section}{Liste e sequenze}

% --- Liste ---
\begin{refbox}[Liste]{csGreen}
\label{sec:liste}

Una lista è una catena di \textit{cons-cell}. Ogni cella contiene
un valore (\fn{car}) e il puntatore al resto (\fn{cdr}).
\fn{NIL} termina la lista.

\begin{lstlisting}
'(a b c)           ; lista letterale (= (cons 'a (cons 'b (cons 'c nil))))
(cons 'a '(b c))   ; -> (A B C)
(car '(a b c))     ; -> A
(cdr '(a b c))     ; -> (B C)
\end{lstlisting}

\textbf{Coppia puntata} (cons non terminante in NIL):\nopagebreak
\begin{lstlisting}
(cons 'a 'b)       ; -> (A . B)  lista impropria
\end{lstlisting}

\textbf{Accesso:}\nopagebreak
\begin{lstlisting}
(first lst)  (rest lst)             ; = car, cdr
(second lst) (third lst)            ; cadr, caddr
(nth 2 lst)                         ; terzo elemento (0-based)
(nthcdr 2 lst)                      ; cdr applicato 2 volte
(last lst)                          ; ultima cons-cell (non elemento!)
(butlast lst)                       ; tutti tranne l'ultimo
\end{lstlisting}

\textbf{Costruzione e modifica:}\nopagebreak
\begin{lstlisting}
(list 'a 'b 'c)             ; -> (A B C)
(append '(a b) '(c d))      ; -> (A B C D)  (non distruttivo)
(reverse '(a b c))          ; -> (C B A)    (non distruttivo)
(nreverse (list 'a 'b 'c))  ; -> (C B A)    (distruttivo!)
(copy-list lst)             ; copia superficiale
\end{lstlisting}

\textbf{Push e pop:}\nopagebreak
\begin{lstlisting}
(push 'x lst)    ; aggiunge in testa (modifica lst)
(pop lst)        ; rimuove e restituisce la testa
\end{lstlisting}

\textbf{Ricerca e test:}\nopagebreak
\begin{lstlisting}
(length '(a b c))           ; -> 3
(member 'b '(a b c))        ; -> (B C)  (non solo T!)
(member 'b '(a b c) :test #'equal)  ; con comparatore custom
(find 'b '(a b c))          ; -> B
(position 'b '(a b c))      ; -> 1
\end{lstlisting}

\textbf{Sostituzione:}\nopagebreak
\begin{lstlisting}
(subst 'x 'b '(a b (b c)))  ; -> (A X (X C))  ricorsivo
(remove 'b '(a b b c))      ; -> (A C)
(remove-duplicates '(a b a c b)) ; -> (A C B)
\end{lstlisting}

\end{refbox}

% --- Predicati ---
\begin{refbox}[Predicati]{csGreen}
\label{sec:predicati}

\textbf{Predicati di tipo:}\nopagebreak
\begin{lstlisting}
(null '())       ; -> T   (lista vuota o NIL)
(atom 'a)        ; -> T   (atomo = non-cons)
(consp '(a))     ; -> T   (cons non vuota)
(listp '(a))     ; -> T   (lista, anche NIL)
(numberp 3)      ; -> T
(integerp 3)     ; -> T
(floatp 3.0)     ; -> T
(symbolp 'foo)   ; -> T
(stringp "x")    ; -> T
(functionp #'+)  ; -> T
(characterp #\a) ; -> T
\end{lstlisting}

\textbf{Uguaglianza:}\nopagebreak
\begin{lstlisting}
(eq 'a 'a)              ; T -- identita' oggetto (simboli, T, NIL)
(eql 3 3)               ; T -- eq + numeri e caratteri
(equal '(a b) '(a b))   ; T -- strutturale (ricorsiva)
(equalp "Ciao" "ciao")  ; T -- strutturale + case-insensitive
\end{lstlisting}

\begin{refnote}[Quando usare quale uguaglianza]
\begin{tabular}{@{}ll@{}}
Numeri:         & \fn{=} \\
Simboli:        & \fn{eq} \\
Caratteri:      & \fn{eql} \\
Liste/strutture: & \fn{equal} \\
Stringhe (no case): & \fn{equalp} o \fn{string-equal} \\
\end{tabular}
\end{refnote}

\end{refbox}

% --- Sequenze ---
\begin{refbox}[Sequenze (liste, vettori, stringhe)]{csGreen}
\label{sec:sequenze}

Molte funzioni funzionano su qualsiasi sequenza:

\begin{lstlisting}
(elt '(a b c) 1)             ; -> B  (lista)
(elt "hello" 1)              ; -> #\e (stringa)
(elt #(10 20 30) 1)          ; -> 20 (vettore)

(length "hello")             ; -> 5
(subseq "hello" 1 3)         ; -> "el"
(concatenate 'list '(a) '(b)) ; -> (A B)
(concatenate 'string "fo" "o") ; -> "foo"
(position #\l "hello")       ; -> 2
(count #\l "hello")          ; -> 2
(remove #\l "hello")         ; -> "heo"

; coerce converte tra tipi di sequenza
(coerce "abc" 'list)         ; -> (#\a #\b #\c)
(coerce '(#\a #\b) 'string)  ; -> "ab"
(coerce #(1 2 3) 'list)      ; -> (1 2 3)
\end{lstlisting}

\end{refbox}

% --- Destructuring ---
\begin{refbox}[Destructuring]{csGreen}
\label{sec:destructuring}

\fn{destructuring-bind} lega variabili alla struttura di una lista:

\begin{lstlisting}
(destructuring-bind (a b c) '(1 2 3)
  (+ a b c))             ; -> 6

(destructuring-bind (head . tail) '(1 2 3)
  tail)                  ; -> (2 3)

(destructuring-bind (a (b c) d) '(1 (2 3) 4)
  (list a b c d))        ; -> (1 2 3 4)
\end{lstlisting}

\end{refbox}

% =============================================================
% MACRO-AREA 4: STRUTTURE DATI
% =============================================================
\addcontentsline{toc}{section}{Strutture dati}

% --- Alist e Plist ---
\begin{refbox}[Association list e Property list]{csTeal}

\textbf{Association list (alist):} lista di coppie \fn{(chiave~.~valore)}:
\label{sec:alist}
\begin{lstlisting}
(defparameter *rubrica*
  '((mario . "333-1111")
    (luigi . "333-2222")))

(assoc 'mario *rubrica*)          ; -> (MARIO . "333-1111")
(cdr (assoc 'mario *rubrica*))    ; -> "333-1111"
(rassoc "333-1111" *rubrica* :test #'equal) ; cerca per valore
(cons '(anna . "333-3333") *rubrica*)  ; aggiunge coppia in testa
(acons 'anna "333-3333" *rubrica*)    ; equivalente, più leggibile
\end{lstlisting}

\textbf{Property list (plist):} chiavi e valori alternati nella stessa lista:\nopagebreak
\begin{lstlisting}
(defparameter *p* '(:nome "Mario" :eta 23))
(getf *p* :nome)           ; -> "Mario"
(getf *p* :via "n/d")     ; -> "n/d"  (default se assente)
(setf (getf *p* :eta) 24)  ; modifica
\end{lstlisting}

\begin{refnote}[Alist vs Plist]
\begin{tabular}{@{}lll@{}}
\textbf{} & \textbf{Alist} & \textbf{Plist} \\
Struttura & \fn{((k . v)...)} & \fn{(k v k v...)} \\
Ricerca   & \fn{assoc}  & \fn{getf} \\
Uso tipico & dati associativi & proprietà oggetti \\
\end{tabular}
\end{refnote}

\end{refbox}

% --- Hash table ---
\begin{refbox}[Hash table]{csTeal}
\label{sec:hashtable}

\begin{lstlisting}
; creazione (:test specifica il comparatore di chiavi)
(make-hash-table)                    ; default: eql
(make-hash-table :test #'equal)      ; chiavi stringa/lista

; lettura (due valori: valore e flag "trovato?")
(gethash 'foo ht)                    ; -> NIL, NIL  (assente)
(gethash 'foo ht "default")          ; -> "default", NIL

; scrittura e rimozione
(setf (gethash 'foo ht) 42)
(remhash 'foo ht)

; utilità
(hash-table-count ht)                ; numero di chiavi presenti
(clrhash ht)                         ; svuota la tabella

; iterazione
(maphash (lambda (k v)
           (format t "~A -> ~A~%" k v))
         ht)
\end{lstlisting}

\begin{refnote}[Ricevere entrambi i valori di gethash]
\begin{lstlisting}
(multiple-value-bind (valore trovato)
    (gethash 'foo ht)
  (if trovato
      (format t "Trovato: ~A~%" valore)
      (format t "Assente~%")))
\end{lstlisting}
\end{refnote}

\end{refbox}

% --- Array ---
\begin{refbox}[Array e vettori]{csTeal}
\label{sec:array}

\begin{lstlisting}
#(10 20 30)          ; vettore letterale (immutabile)

; creazione (sempre con :initial-element per evitare letture indefinite)
(make-array 5 :initial-element 0)           ; vettore di 5 zeri
(make-array 5 :initial-element nil)         ; vettore di 5 NIL
(make-array '(3 3) :initial-element 0)      ; matrice 3x3

; accesso e modifica
(aref v 2)                                  ; terzo elemento
(aref m 1 2)                                ; riga 1, colonna 2
(setf (aref v 0) 99)

; dimensioni
(length v)                                  ; lunghezza (vettori 1D)
(array-dimensions m)                        ; -> (3 3)
(array-dimension m 0)                       ; -> 3 (numero righe)

; vettore dinamico (adjustable)
(defparameter *v*
  (make-array 0 :fill-pointer 0 :adjustable t))
(vector-push-extend 42 *v*)                 ; aggiunge in coda
(vector-pop *v*)                            ; rimuove dall'ultima posizione
\end{lstlisting}

\end{refbox}

% --- Defstruct ---
\begin{refbox}[Strutture con defstruct]{csTeal}
\label{sec:strutture}

\fn{defstruct} genera automaticamente costruttore, accessori,
predicato e funzione di copia:

\begin{lstlisting}
(defstruct persona
  nome
  (eta 18)        ; slot con valore di default
  citta)

; costruzione
(defparameter *p*
  (make-persona :nome "Mario" :eta 23 :citta "Roma"))

; accesso agli slot
(persona-nome *p*)           ; -> "Mario"
(persona-eta *p*)            ; -> 23

; modifica
(setf (persona-eta *p*) 24)

; predicato e copia
(persona-p *p*)              ; -> T
(copy-persona *p*)           ; copia superficiale
\end{lstlisting}

\begin{refnote}[Cosa genera defstruct automaticamente]
\begin{tabular}{@{}ll@{}}
\fn{make-nome}    & costruttore \\
\fn{nome-slot}    & lettore per ogni slot \\
\fn{(setf (nome-slot ...) ...)} & modificatore \\
\fn{nome-p}       & predicato di tipo \\
\fn{copy-nome}    & copia superficiale \\
\end{tabular}
\end{refnote}

\end{refbox}

% --- Multiple values ---
\begin{refbox}[Valori multipli]{csTeal}
\label{sec:multiplevalues}

Una funzione può restituire più di un valore con \fn{values}:

\begin{lstlisting}
(values 3 4)                ; restituisce due valori

; cattura con multiple-value-bind
(multiple-value-bind (q r)
    (floor 17 5)
  (format t "~A resto ~A~%" q r))  ; "3 resto 2"

; raccoglie i valori in una lista
(multiple-value-list (floor 17 5)) ; -> (3 2)

; passa i valori multipli come argomenti
(multiple-value-call #'+ (values 1 2) (values 3 4)) ; -> 10
\end{lstlisting}

\end{refbox}

% =============================================================
% MACRO-AREA 5: ASTRAZIONE
% =============================================================
\addcontentsline{toc}{section}{Funzioni, closures e macro}

% --- Higher-order ---
\begin{refbox}[Funzioni higher-order]{csPurple}
\label{sec:higherorder}

\fn{\#'} ottiene l'oggetto funzione; \fn{lambda} crea funzioni anonime:

\begin{lstlisting}
#'quadrato           ; oggetto funzione (= (function quadrato))
(lambda (x) (* x x)) ; funzione anonima

(funcall #'quadrato 5)         ; -> 25
(funcall (lambda (x) (* x x)) 5) ; -> 25
(apply #'+ '(1 2 3))           ; -> 6
(apply #'max 3 '(1 5 2))       ; -> 5  (argomenti misti)
\end{lstlisting}

\textbf{Funzioni sulle liste:}\nopagebreak
\begin{lstlisting}
(mapcar #'quadrato '(1 2 3 4))          ; -> (1 4 9 16)
(mapcar #'+ '(1 2 3) '(10 20 30))       ; -> (11 22 33)
(mapcan #'list '(1 2 3))                ; -> (1 2 3)
(remove-if #'evenp '(1 2 3 4 5))        ; -> (1 3 5)
(remove-if-not #'oddp '(1 2 3 4 5))     ; -> (1 3 5)
(find-if #'evenp '(1 3 4 5))            ; -> 4
(find-if-not #'evenp '(1 3 4 5))        ; -> 1
(every #'numberp '(1 2 3))              ; -> T
(some  #'null '(1 nil 3))               ; -> T
(count-if #'evenp '(1 2 3 4))           ; -> 2
(reduce #'+ '(1 2 3 4))                 ; -> 10
(reduce #'+ '() :initial-value 0)       ; -> 0  (lista vuota)
(sort (copy-list '(3 1 2)) #'<)         ; -> (1 2 3)  (distruttivo!)
(stable-sort lst #'<)                   ; preserva ordine degli uguali
\end{lstlisting}

\begin{refnote}[sort è distruttivo]
\fn{sort} può modificare la lista originale. Usare sempre
\fn{(sort (copy-list lst) \#'<)} se si vuole preservare \fn{lst}.
\end{refnote}

\end{refbox}

% --- Closures ---
\begin{refbox}[Closures]{csPurple}
\label{sec:closures}

Una \textit{closure} è una funzione che cattura le variabili del
proprio ambiente lessicale. È il meccanismo principale per
incapsulare stato privato:

\begin{lstlisting}
(defun crea-contatore ()
  (let ((n 0))
    (lambda ()
      (incf n)     ; n e' privata, non accessibile dall'esterno
      n)))

(defparameter *c1* (crea-contatore))
(defparameter *c2* (crea-contatore))  ; stato indipendente da *c1*

(funcall *c1*)   ; -> 1
(funcall *c1*)   ; -> 2
(funcall *c2*)   ; -> 1  (contatore separato)
\end{lstlisting}

\textbf{Closure con più operazioni:}\nopagebreak
\begin{lstlisting}
(defun crea-generatore (start step)
  (let ((curr start))
    (lambda ()
      (prog1 curr        ; ritorna curr, poi lo incrementa
        (incf curr step)))))

(defparameter *gen* (crea-generatore 0 2))
(funcall *gen*)  ; -> 0
(funcall *gen*)  ; -> 2
(funcall *gen*)  ; -> 4
\end{lstlisting}

\end{refbox}

% --- Macro ---
\begin{refbox}[Macro]{csPurple}
\label{sec:macro}

Le macro trasformano codice \textit{prima} della valutazione
(espansione a compile-time). Ricevono le forme non valutate.

\begin{lstlisting}
(defmacro my-when (test &body body)
  `(if ,test
       (progn ,@body)
       nil))

; uso
(my-when (> x 0)
  (print "positivo")
  x)

; verifica dell'espansione
(macroexpand-1 '(my-when (> x 0) (print x)))
; -> (IF (> X 0) (PROGN (PRINT X)) NIL)
\end{lstlisting}

\begin{itemize}[leftmargin=1.5em, itemsep=2pt]
  \item \fn{`} (backquote) --- crea una lista quasi-letterale
  \item \fn{,} (unquote) --- valuta l'espressione nel contesto del chiamante
  \item \fn{,@} (splicing) --- inserisce gli elementi di una lista
\end{itemize}

\textbf{Igiene con \fn{gensym}:} evita la cattura accidentale di variabili:\nopagebreak
\begin{lstlisting}
(defmacro my-swap (a b)
  (let ((tmp (gensym "TMP")))  ; genera simbolo unico
    `(let ((,tmp ,a))
       (setf ,a ,b)
       (setf ,b ,tmp))))
\end{lstlisting}

\end{refbox}

% =============================================================
% MACRO-AREA 6: I/O, ERRORI, CLOS, PATTERN, CONVENZIONI
% =============================================================
\addcontentsline{toc}{section}{I/O e gestione errori}

% --- I/O ---
\begin{refbox}[Input/Output]{csTeal}
\label{sec:io}

\textbf{Lettura:}\nopagebreak
\begin{lstlisting}
(read)             ; legge una S-expression da stdin
(read-line)        ; legge una riga come stringa
(read-char)        ; legge un singolo carattere
\end{lstlisting}

\textbf{Stampa:}\nopagebreak
\begin{lstlisting}
(print x)          ; stampa con quote + newline prima
(princ x)          ; stampa senza quote (leggibile)
(prin1 x)          ; stampa con quote, senza newline prima
(terpri)           ; solo newline
(fresh-line)       ; newline solo se non gia' a inizio riga
\end{lstlisting}

\textbf{format:} funzione di output formattato:\nopagebreak
\begin{lstlisting}
(format t "~A~%"     x)    ; generico + newline
(format t "~D~%"     n)    ; intero decimale
(format t "~,2F~%"   f)    ; float con 2 decimali
(format t "~S~%"     x)    ; con quote (come prin1)
(format t "~{~A ~}~%" lst) ; itera su lista
(format t "~A e ~A~%" a b) ; argomenti multipli

; format su stringa (nil = non stampare, ritorna stringa)
(format nil "Risultato: ~A" val)  ; -> "Risultato: 42"
\end{lstlisting}

\textbf{Direttive format principali:}

\smallskip
\begin{tabular}{@{}l@{\hspace{8pt}}l@{\hspace{16pt}}l@{\hspace{8pt}}l@{}}
\fn{\textasciitilde A} & generico (aesthetic) &
\fn{\textasciitilde S} & con escape \\
\fn{\textasciitilde D} & intero dec &
\fn{\textasciitilde F} & float \\
\fn{\textasciitilde\%} & newline &
\fn{\textasciitilde\&} & fresh-line \\
\fn{\textasciitilde\textasciitilde} & tilde letterale &
\fn{\textasciitilde nT} & tab a colonna n \\
\end{tabular}

\medskip
\textbf{File:}\nopagebreak
\begin{lstlisting}
; lettura riga per riga
(with-open-file (s "dati.txt" :direction :input)
  (loop for riga = (read-line s nil nil)
        while riga
        collect riga))

; scrittura
(with-open-file (s "out.txt"
                   :direction :output
                   :if-exists :supersede)
  (format s "Ciao~%"))
\end{lstlisting}

\end{refbox}

% --- Gestione errori ---
\begin{refbox}[Gestione errori]{csTeal}
\label{sec:errori}

\begin{lstlisting}
; segnalare un errore
(error "Valore non valido: ~A" x)
(warn  "Attenzione: ~A" msg)      ; segnalazione non fatale

; gestire errori con handler-case
(handler-case
    (/ 10 0)
  (division-by-zero (e)
    (format t "Divisione per zero~%"))
  (error (e)
    (format t "Errore: ~A~%" e)))

; ignorare errori (ritorna NIL + condizione)
(ignore-errors
  (/ 10 0))    ; -> NIL, #<DIVISION-BY-ZERO ...>

; handler-bind (non unwind lo stack)
(handler-bind
    ((error (lambda (e)
              (format t "Intercettato: ~A~%" e))))
  (/ 10 0))
\end{lstlisting}

\end{refbox}

% =============================================================
\addcontentsline{toc}{section}{CLOS (programmazione a oggetti)}
% =============================================================

\begin{refbox}[CLOS --- Object-Oriented \normalfont\small(argomento avanzato)]{csPurple}
\label{sec:clos}

Il \textbf{Common Lisp Object System} (CLOS) è il sistema di
classi integrato in ANSI Common Lisp. Supporta ereditarietà
multipla e metodi polimorfici.

\begin{lstlisting}
; definizione di classe
(defclass animale ()                       ; () = nessuna superclasse
  ((nome  :initarg :nome  :accessor animale-nome)
   (verso :initarg :verso :reader   animale-verso)))

; sottoclasse
(defclass cane (animale)
  ())

; generic function e metodo
(defgeneric descrivi (a))

(defmethod descrivi ((a animale))
  (format t "~A dice ~A~%"
          (animale-nome a)
          (animale-verso a)))

; specializzazione nella sottoclasse
(defmethod descrivi ((c cane))
  (format t "Il cane ~A abbaia!~%" (animale-nome c)))

; creazione istanza
(defparameter *fido*
  (make-instance 'cane :nome "Fido" :verso "Bau"))

(descrivi *fido*)       ; chiama il metodo specializzato per CANE

; predicato di tipo
(typep *fido* 'cane)    ; -> T
(typep *fido* 'animale) ; -> T  (ereditarieta')
\end{lstlisting}

\fn{defmethod} può anche definire metodi \fn{:before},
\fn{:after}, \fn{:around} per decorare il comportamento
senza sovrascrivere il metodo principale.

\end{refbox}

% =============================================================
\addcontentsline{toc}{section}{Pattern e convenzioni}
% =============================================================

\begin{refbox}[Pattern essenziali]{csOrange}
\label{sec:pattern}

\textbf{Elaborazione ricorsiva di lista:}\nopagebreak
\begin{lstlisting}
(defun trasforma-lista (lst)
  (if (null lst)
      '()
      (cons (trasforma (car lst))
            (trasforma-lista (cdr lst)))))
\end{lstlisting}

\textbf{Accumulatore tail-recursive:}\nopagebreak
\begin{lstlisting}
(defun somma (lst)
  (labels ((iter (lst acc)
             (if (null lst)
                 acc
                 (iter (cdr lst) (+ acc (car lst))))))
    (iter lst 0)))
\end{lstlisting}

\textbf{Memoization:}\nopagebreak
\begin{lstlisting}
(defun memoize (f)
  (let ((cache (make-hash-table :test #'equal)))
    (lambda (&rest args)
      (multiple-value-bind (val trovato)
          (gethash args cache)
        (if trovato val
            (setf (gethash args cache)
                  (apply f args)))))))
\end{lstlisting}

\textbf{Backtracking:}\nopagebreak
\begin{lstlisting}
(defun risolvi (stato)
  (cond
    ((soluzione-p stato) (list stato))
    ((vicolo-cieco-p stato) nil)
    (t (loop for mossa in (mosse-valide stato)
             append (risolvi (applica mossa stato))))))
\end{lstlisting}

\end{refbox}

\begin{refbox}[Convenzioni di stile]{csBlue}
\label{sec:stile}

\begin{itemize}[leftmargin=1.5em, itemsep=3pt]
  \item Nomi con \textbf{kebab-case}: \fn{mia-funzione}, \fn{conta-elementi}
  \item Variabili globali dinamiche con \textbf{earmuffs}: \fn{*contatore*}
  \item Costanti con \textbf{più segni}: \fn{+max-dimensione+}
  \item Predicati terminano in \textbf{-p}: \fn{null}, \fn{evenp}, \fn{palindromo-p}
  \item Reimplementazioni di built-in: suffisso \fn{*}: \fn{length*}, \fn{reverse*}
  \item Ogni funzione pubblica dovrebbe avere una \textit{docstring}
  \item Funzioni piccole e composabili; evitare stato globale
  \item Gestire esplicitamente i casi limite: lista vuota, \fn{nil}, zero
\end{itemize}

\end{refbox}

\begin{refbox}[Package]{csBlue}
\label{sec:package}

\begin{lstlisting}
(defpackage :mio-progetto
  (:use :cl)
  (:export #:funzione-pubblica
           #:altra-funzione))

(in-package :mio-progetto)
\end{lstlisting}

Durante la fase di apprendimento è sufficiente lavorare nel
package \fn{cl-user} (il default al REPL).
Per esercizi multi-file, definire un package per progetto
ed esportare solo le funzioni pubbliche.

\end{refbox}
